\documentclass[11pt]{article}
\usepackage[margin=1in]{geometry}
\usepackage[T1]{fontenc}
\usepackage[utf8]{inputenc}
\usepackage{lmodern}
\usepackage{amsmath,amssymb}
\usepackage[hidelinks]{hyperref}
\title{Why Space and Time Must Change: The Invariance of $c$}
\author{John C. W. McKinley\\Independent Researcher\\ORCID: 0009-0005-7097-5035}
\date{November 20, 2026}
\begin{document}
\maketitle
\noindent\footnotesize\textit{Working draft published at \href{https://johncwmckinley.com}{johncwmckinley.com}. Not deposited on Zenodo; subject to revision.}\normalsize
\noindent\textbf{Status:} Theoretical interpretation --- working manuscript

\begin{abstract}
Space is not an independent container. In the relational ontology of the Timeless Light Model, physical reality consists of lawful relationships and interactions, while registrations and retained records support the reconstruction of spatial configurations and histories. The familiar picture of a universe that must morph to honor the invariant speed $c$ is a metaphor: no separate spatial stage is resized. Special relativity already supplies a continuous, coordinated dependence of spatial and temporal measurements on relative velocity, including relativity of simultaneity, while preserving the same causal structure. A particle circulating through an accelerator illustrates the distinction between different accumulated proper times and participation in one set of lawful physical interactions. The contribution is interpretive; no established equations or predictions are changed.
\end{abstract}
\noindent\textbf{Keywords:} relational space, invariant causal speed, Lorentz transformations, registration, proper time, Timeless Light Model

\section{The Question Behind the Metaphor}

The universe has to morph to honor the invariant speed limit $c$.

The sentence supplies an intuitive picture, not a proposed physical mechanism. An observer moving relative to another assigns different distances and time coordinates to events, yet both descriptions preserve the invariant causal-speed relationship. The familiar picture invites the thought that a completed universe must adjust its dimensions whenever an observer's motion changes.

That picture assumes there is an independently existing space to adjust.

The latest relational formulation of the Timeless Light Model (TLM) rejects that assumption. Space is not a location or a container for physical relationships. It is the geometric description of relationships among physical systems and events \cite{mckinleylocation2026,mckinleybox2026}. Motion is a real, measurable change of relationships, not proof of displacement through an additional spatial substance \cite{mckinleymotion2026}.

The explanatory question therefore changes. It is no longer: \emph{What physically reshapes the universe to accommodate an observer?} It becomes: \emph{Why do different relational measurements remain consistent with the same invariant causal constraint?}

Special relativity already answers the mathematical part. The present paper addresses the interpretation when space is not presumed to be an independent place.

\section{Lawful Relationships, Registration, and Records}

TLM begins with timeless governing laws. Such laws specify admissible relationships, interactions, and outcomes; they are not themselves events occurring inside a larger spatial or temporal container \cite{mckinleybox2026}. Physical actuality consists of lawful configurations and interactions, independently of whether a particular system has acquired information about them.

A \emph{registration} is a local physical outcome of a lawful interaction. A detector response, an instrumental measurement, or an encounter between physical systems supplies an example. Registration does not require a person to observe it; the lawful relationship does not depend on being measured.

A \emph{record} is distinguishable from the registration it preserves. A clock reading, detector trace, or stored measurement makes information about a registration available for later comparison. A system evaluating multiple records reconstructs separations, shapes, trajectories, durations, and constraints. Those reconstructions are anchored in actual interactions and remain subject to invariant physical laws.

The distinction matters because an observed spatial world is not evidence for a second physical entity called the spatial stage. Registered relationships support a coherent spatial description; they do not thereby establish a completed container in which the relationships were housed.

\section{The Flipbook: Records Without a Persisting Past}

Consider a flipbook showing a donkey standing beside a clown, raising its hind leg, extending that leg, and the clown recoiling. The reader encounters the drawings successively. At each encounter, only the current observation and the reader's presently retained memories of earlier observations are available. Those earlier moments of reading have passed; they do not continue to exist as separate physical nows.

Comparing the successive registrations, the reader deduces: \emph{the donkey kicked the clown}. The deduction rests on a remembered sequence of recorded positions and changes. It also supports a prediction: the reader expects a later page to show the clown on the floor. The record thus serves two purposes: reconstructing what occurred and anticipating what another registration will reveal.

The flipbook is only an illustration. The point does not depend on gaps between drawings, on whether an animation is continuous, or on denying the reality of physical movement. A camera, detector, memory, or any other retained record supplies information about earlier physical conditions from which a reader can infer relationships and make predictions. Those inferences are constrained by lawful physical interactions, not created by the reader.

In TLM's universal-present interpretation, \emph{there is only now}. Previous nows have no continued physical existence. What exists now includes the physical records and dynamical conditions left by earlier events; those present conditions support deductions about the past and predictions about later registrations. Neither the deduction nor the prediction requires past or future configurations to remain physically present \cite{mckinleybox2026,mckinleymotion2026}.

\section{The Continuous Relativistic Relationship}

Special relativity establishes that the vacuum speed of light has the same value in every inertial reference frame, alongside the principle of relativity \cite{einstein1905,taylorwheeler1992}. The Lorentz transformations give the coordinated assignment of spatial and temporal coordinates that preserves this invariant relationship.

Those assignments vary continuously with relative velocity. They do not switch on at a special speed. As velocity changes, the measured longitudinal length of a moving system, the comparison of clock rates, and the assignment of simultaneity vary according to one connected relativistic geometry.

The familiar Lorentz factor expresses the nonlinear dependence on relative speed $v$:
\[
\gamma(v)=\frac{1}{\sqrt{1-v^2/c^2}},\qquad |v|<c.
\]
At ordinary speeds the resulting differences are small. As the magnitude of relative velocity approaches $c$, $\gamma$ grows without bound and its reciprocal approaches zero. These are limiting features of established mathematical relationships, not physical adjustments performed by a separate universe.

The variables do not add up to $c$, and length contraction alone does not establish invariance. The \emph{complete} Lorentz transformation, including relativity of simultaneity, preserves the spacetime interval and light-cone structure \cite{minkowski1909,taylorwheeler1992}. Frame-dependent assignments differ; the invariant constraint does not.

Proper time is an essential distinction. A particular clock's accumulated proper time along its worldline is not revised when another observer adopts a different reference frame. Different physical paths accumulate different proper times. This is a physical fact about those paths, not a disagreement caused by observing them.

No rederivation of the Lorentz transformations is needed here. Their established, smoothly velocity-dependent relationships supply the precise consistency that the metaphor of a morphing universe was intended to convey.

\section{The Collider: Different Proper Times, One Physical Relationship}

Consider a proton circulating at a speed close to $c$ through a beam pipe in the Large Hadron Collider. The evacuated beam pipe and accelerator structures remain approximately at rest in the laboratory frame; electromagnetic fields accelerate, guide, and focus the beam \cite{cernlhc,cernaccelerator}.

The proton follows a timelike path that accumulates far less proper time, over a given laboratory interval, than equipment fixed to the accelerator. Yet the proton remains in lawful relation to that equipment. Its trajectory is constrained by accelerator fields, and passages and interactions are registered by instrumentation. The apparatus need not share the proton's accumulated proper time in order to interact with it.

An instructive comparison uses two passages of an idealized circulating particle past the same fixed beamline position. A laboratory clock measures the elapsed interval between passages, while the particle accumulates proper time along its intervening curved path. The elapsed amounts differ according to relativity. Their continued physical proximity at the ring presents no contradiction.

This is not a case of the particle existing in one temporal world and the beam pipe in another. Proper-time totals belong to different physical histories. Their interactions and the records of those interactions belong to one consistent causal structure.

In the relational interpretation, the question of how they remain close while accumulating different amounts of time begins from the wrong image. There is no independently existing place called space in which incompatible times must be reconciled. There are lawful spatial and temporal relationships, locally accumulated proper times, and physical interactions consistent with the invariant causal constraint.

The particle in a circular accelerator is accelerated; it has no single inertial rest frame covering its entire orbit. The proper-time comparison therefore concerns its worldline, not one comoving Lorentz frame applied globally around the ring.

The collider example is a consistency demonstration, not an independent experimental proof of the relational ontology. It makes the explanatory content concrete: unequal proper-time histories do not require separate physical actualities.

\section{The Present Is Not a Universal Clock}

TLM's recent relational account proposes one physical actuality, rather than separately existing past, present, and future spatial containers \cite{mckinleybox2026,mckinleymotion2026}. Each physical system participates in lawful relationships and accumulates proper time locally along its history.

That ontological present must not be confused with a universal simultaneity coordinate. Special relativity permits different inertial frames to assign different times to spatially separated events. This interpretation does not select a preferred reference frame or provide instantaneous knowledge of distant events.

Information becomes available through causally constrained interactions and registration. Records permit later comparison. A system reconstructs earlier events from present evidence without requiring a physically persisting past, and assigns spatial descriptions without positing an independent spatial background.

The invariant $c$ constrains which events can be causally connected and how their spatial and temporal assignments fit together. The orderly appearance reconstructed from registrations therefore expresses lawful relationships constrained from the outset.

\section{Interpretive Result}

The argument is an explanatory ordering, not a replacement for special relativity:
\begin{quote}
\emph{Timeless lawful structure governs physical relationships; local interactions register outcomes; retained records support reconstructions of spatial and temporal histories; and the Lorentz relationships preserve the same invariant causal structure across reference frames.}
\end{quote}

The flipbook clarifies how present records of earlier registrations support deductions about movement and predictions of later outcomes, without a physically persisting past. The collider clarifies why unequal proper-time histories do not require distinct temporal worlds. The established Lorentz relationships show why measurements vary continuously with relative velocity without violating invariant $c$.

Nothing here changes the formulas or predictions of relativity. The proposed change is to the ontology assigned to their success: relational measurements do not require a universe-container that must physically morph to accommodate them.

\section{Conclusion}

The invariant speed $c$ is a fundamental constraint governing lawful physical relationships. Space is not an independent container that must change to preserve it. Spatial and temporal measurements express relationships that are already consistent with the invariant causal structure.

Different reference frames assign different distances, coordinate times, and relations of simultaneity. These assignments vary continuously and coherently according to the Lorentz transformations. Length contraction and time dilation are expressions of this same lawful structure, not separate physical adjustments performed by the universe.

Physical interactions produce registrations, and retained records permit the reconstruction of spatial configurations, motion, and elapsed time. The resulting appearance of a universe changing through space and time does not require an independently existing spatial stage or a physically persisting past. Different systems accumulate their own proper times while participating in the same lawful physical reality.

\textbf{Nothing needs to morph to preserve $c$. The relationships already honor it.}

The universe does not change its laws. Different frames assign different measurements within the same lawful order.

\begin{thebibliography}{10}
\bibitem{einstein1905} Einstein, A. (1905). \emph{On the Electrodynamics of Moving Bodies.} Annalen der Physik, 17, 891--921.
\bibitem{minkowski1909} Minkowski, H. (1909). \emph{Space and Time.} Physikalische Zeitschrift, 10, 104--111.
\bibitem{taylorwheeler1992} Taylor, E. F., and Wheeler, J. A. (1992). \emph{Spacetime Physics}, 2nd ed. W. H. Freeman.
\bibitem{mckinleylocation2026} McKinley, J. C. W. (28 October 2026). \emph{Space Is Not a Location. It Is a Relationship: A Relational Ontology of Distance Without an Enclosing Universe.} Website working draft. \url{https://johncwmckinley.com/papers/space-is-not-a-location/}
\bibitem{mckinleybox2026} McKinley, J. C. W. (28 October 2026). \emph{Space Is Not a Box: Timeless Law, Universal Actuality, and the Registration of Space and Motion.} Website working draft. \url{https://johncwmckinley.com/papers/space-is-not-a-box/}
\bibitem{mckinleymotion2026} McKinley, J. C. W. (18 October 2026). \emph{Motion Is Relational: A No-Go on Independent Position and Substance Displacement. The Universal Present, Local Accrual, and Lawful Physical Relations.} Website working draft. \url{https://johncwmckinley.com/papers/motion-is-relational/}
\bibitem{cernlhc} CERN. \emph{The Large Hadron Collider.} \url{https://home.cern/science/accelerators/large-hadron-collider/}
\bibitem{cernaccelerator} CERN. \emph{How an Accelerator Works.} \url{https://home.cern/science/accelerators/how-accelerator-works/}
\end{thebibliography}
\end{document}
